%=========================
% Plantilla LaTeX Libro
% Compilar con: pdflatex + biber + pdflatex + pdflatex
%=========================
\documentclass[11pt,a4paper]{book}

%-------------------------
% Idioma y tipografía
%-------------------------
\usepackage[english]{babel}
\usepackage{csquotes}
\usepackage[T1]{fontenc}
\usepackage{lmodern}

%-------------------------
% Márgenes y espaciado
%-------------------------
\usepackage[a4paper,margin=2.5cm]{geometry}
\usepackage{setspace}
\onehalfspacing

%-------------------------
% Matemáticas y símbolos
%-------------------------
\usepackage{amsmath,amssymb,amsthm}

%-------------------------
% Tablas y figuras
%-------------------------
\usepackage{graphicx}
\usepackage{booktabs}
\usepackage{caption}
\usepackage{subcaption}
\usepackage{changepage}

%-------------------------
% Enlaces y PDF
%-------------------------
\usepackage{hyperref}
\hypersetup{
  colorlinks=true,
  linkcolor=blue,
  citecolor=blue,
  urlcolor=blue,
  pdftitle={On the indefiniteness of the continuum hypothesis},
  pdfauthor={Rodrigo Blanco Betes}
}

%-------------------------
% Listas, código y utilidades
%-------------------------
\usepackage{enumitem}
\usepackage{xcolor}
\usepackage{csquotes} % recomendado con biblatex

%-------------------------
% Bibliografía (biblatex + biber)
%-------------------------
\usepackage[
  backend=biber,
  style=apa,
  sortcites=true,
  url=true,
  doi=true
]{biblatex}
\DeclareLanguageMapping{english}{english-apa}
\addbibresource{refs.bib}

%-------------------------
% Metadatos del artículo
%-------------------------
\title{\textbf{On the indefiniteness of the continuum hypothesis}}

\author{Rodrigo Blanco Betes}


\date{\today}

%-------------------------
% Entornos de teoremas (opcional)
%-------------------------
\newtheorem{theorem}{Theorem}
\newtheorem{lemma}{Lemma}

\begin{document}

% ---------- Portada ----------
\frontmatter
\maketitle

\tableofcontents

%-------------Introducción
\section*{Introduction}


%----------------Texto principal----------
\mainmatter

\chapter{Background}

%-------------Primeros antecedentes: introducción lógica a la hipótesis del continuo
\section{The Development of CH: From a Mathematical Problem to a Philosophical Question}

%-------------Versiones de filosofía de las matemáticas
\section{Philosophy of mathematics in the early twentieth century: the big three}
Let us recall the liar paradox, the statement that “This statement is false”. \textcite{Russell1908} rephrased it as follows: “It is not true for all propositions $p$ that if I affirm $p$, $p$ is true”. He then concluded that the cause of the paradox was a quantification over all the propositions, so that the notion of an arbitrary proposition is illegitimate. A similar reasoning was employed to explain the barber paradox, which appears after trying to define the naïve set ${x \mid x \notin x}$. This concept of an arbitrary naïve set was considered to be unclear, non-definite: this is what was thought to lead to the contradiction. Intuitionists tried to extend this debate even further, by stating that these paradoxes come from accepting logical structures that do not have a counterpart in our intuition (such as the law of excluded middle). We could perhaps say something similar about the concept of an arbitrary subset of the continuum: maybe this concept is so far from our intuition that trying to understand it in linguistic terms may lead us to contradictions or at least incongruities. However, there is no compelling evidence that the concept of an arbitrary subset of the continuum can lead to a contradiction in the same sense as the concept of a naïve set did, or the concept of all propositions. In fact, axiom systems such as ZF are hardly ever doubted to be consistent.

%------------Lógica vs matemáticas. Por qué las matemáticas son sintéticas
\section{Our stand towards logicism and the meaning of mathematical statements}

In addition to this, the idea of clarity of a determinate concept seems to be too subjective to constitute a criterion for its validity. This is something Koellner also points out:

\begin{adjustwidth}{1cm}{1cm}
\begin{quote}
\small
Intuitions of clarity are not robust. Some people think that the concept of an arbitrary subset of natural numbers is about as clear a conception as one can have; others disagree. Whose intuitions are to count? \textcite[p. 14]{Koellner2022}.
\end{quote}
\end{adjustwidth}

\begin{adjustwidth}{1cm}{1cm}
\begin{quote}
\small
There is no problem to put oneself in the mental frame of mind of “this is what the cumulative hierarchy looks like”, for which one can see that such and such propositions including the axioms of ZFC are (more or less) obviously true. I have taught set theory many times and have presented it in terms of this ideal-world picture with only the caveat that this is what things are supposed to be like in that world, rather than to assert that that's the way the world actually is. \textcite[p. 19]{Feferman2011a}.
\end{quote}
\end{adjustwidth}

Feferman's disbelief in a Platonic set-theoretical realm is again evident. His skepticism, however, does not seem to stem from a difficulty in conceiving the realm, but from the distance between such idealized structures and the real world. He concludes the paragraph with a comparison to idealized models in physics:

\begin{adjustwidth}{1cm}{1cm}
\begin{quote}
\small
I am no more uncomfortable doing that than if I were teaching Newtonian mechanics and talking of a world in which space-time is Euclidean (...), or if I were teaching relativistic mechanics. And one can see why little more is needed conceptually to develop an extraordinary part of set theory; as in the case of number theory, a little bit goes a long way. The issue is not whether we can think in such terms but, rather, what the philosophical status is of that kind of thought. \textcite[p. 19]{Feferman2011a}.
\end{quote}
\end{adjustwidth}



\end{document}